% STATUS: reviewed
\chapter{Sistema ir kursas}\label{ch:course}

\product{} buvo sukurtas kaip kurso projektas, o kurso metodas yra jo specifikacija. Šiame
skyriuje abu palyginami greta: ko prašo trys kurso užduotys ir kur produkte yra kiekvienos
atsakymas; kurios kurso sąvokos tapo kuriomis sistemos dalimis; kur sistema seka paties šaltinio
tekstą, o ne jo santrauką; ir ko kursas moko, bet sistema nedaro.

\section{Kursas ir jo pavyzdinė organizacija}\label{sec:course:course}

Kursas --- \enquote{Finansų teisė ir finansų įstaigų veiklos atitikties valdymas} --- dėstomas
Vilniaus universiteto Kauno fakulteto finansinių technologijų (FinTech) magistrantūros
programoje 2026 m. rudenį. Jo paskaitos apima finansų rinkų struktūrą ir kiekvienos rūšies
įstaigų licencijavimą, atitikties funkcijos vietą vidaus valdyme ir atitikties valdymo sistemos
kūrimą, pinigų plovimo prevencijos ir skaitmeninio veiklos atsparumo reikalavimus, taip pat
atsakomybę už pažeidimus. Savarankiškas darbas atliekamas su viena pavyzdine organizacija,
\demo{}: tai imituotas kriptoturto paslaugų teikėjas, kuriame dirba apie keturiasdešimt žmonių,
paskelbtas pagal \en{Creative Commons} licenciją su 26 politikomis --- trylika informacijos
saugumo valdymo sistemos ir trylika jo finansinių technologijų verslo politikų
(\cref{ch:company}).

Kursas taip pat aprašo rolę: technologą, kuriantį įrankius teisiniam ir atitikties darbui. Šiai
rolei priskirtos pareigos skamba kaip tokios sistemos kontrolinis sąrašas --- instrukcijos,
kurių atsakymai susieti su schema, atsakymai, pagrįsti surastais šaltiniais, regresinių testų
rinkiniai, kuriuose vienas modelis vertina kitą, saugomi etaloniniai atvejų rinkiniai ir audito
sekos. Jas aptaria \cref{sec:course:duties}.

\section{Trys kurso užduotys}\label{sec:course:assignments}

Kiekvienoje kurso užduotyje studentas turi parengti dokumentų rinkinį apie \demo{}, pasitelkdamas
kalbos modelį. Antrąją ir trečiąją užduotį sistema traktuoja kaip produktus: tai, kas sugeneruota
iš pačios įmonės dokumentų, kodu patikrinta pagal užduoties nustatytą formą ir paliekama žmogui
redaguoti ir patvirtinti; į pirmąją ji atsako pokalbyje ir išsaugo rezultatą įmonės duomenyse.
\Cref{tab:course:assignments} pateikia atitiktį; po jos einančios pastraipos paaiškina, kokia
yra forma ir kaip užtikrinama, kad jos būtų laikomasi.

\begin{table}[H]
\caption{Kurso užduotys ir kur produkte yra kiekviena}\label{tab:course:assignments}
\begin{tabularx}{\textwidth}{@{}L{22mm}YY@{}}
\toprule
\hd{Kurso užduotis} & \hd{Ko ji prašo} & \hd{Kur ji yra produkte} \\
\midrule
1, pirmoji dalis & Klasifikuoti organizaciją: finansų rinka, kurioje ji veikia, ją
  reglamentuojantys finansų teisės institutai ir jai reikalinga licencija ar įrašas į sąrašą. &
  Pokalbis su pokalbių asistentu, kuris atsako iš teisės aktų bazės su nuorodomis ir savo
  įrankiais įrašo rezultatą į įmonės profilį (\code{financial_market})
  (\cref{ch:chat,ch:company}). \\
1, antroji dalis & Rolės profilis --- pavyzdžiui, atitikties pareigūno --- parengtas pagal
  pačios įmonės dokumentus. &
  Įmonės rolė (\code{company_roles}); pokalbis gali ją pasirinkti, ir tada jos tekstas patenka
  į pokalbio instrukciją kaip kintamasis \code{ROLE_CONTEXT} (\cref{ch:chat}). \\
2 & Atitikties politika iš aštuonių dalių ir trūkstamų antrojo lygmens politikų bei trečiojo
  lygmens procedūrų sąrašai; toliau --- dokumentuose rastos pareigybės, organizacinė
  struktūra, atitikties procesas etapais ir atsakomybių matrica. &
  Programos dokumentai \code{policy} ir \code{governance}, eksportuojami PDF formatu, o
  politika --- ir Word formatu (\cref{ch:programme}). \\
3 & Tikimybės ir poveikio kriterijai, $5 \times 5$ matrica su rizikos lygiais ir atsako į riziką būdais, ne
  mažiau kaip penki rizikos scenarijai nustatyto pavidalo sakiniu ir teisinės rizikos registras. &
  Programos dokumentas \code{risks}, eksportuojamas PDF ir Excel formatais
  (\cref{ch:programme}). \\
\bottomrule
\end{tabularx}
\end{table}

\paragraph{2 kurso užduotis: politika.} Kursas moko atitikties politikos iš aštuonių dalių, ir
produkto politika turi lygiai šias aštuonias dalis ta pačia tvarka: įvadą su tikslu, taikymo
sritimi ir sąsaja su kitais dokumentais; vadovybės įsipareigojimą ir jo ryšį su įmonės
vertybėmis; atitikties funkcijos įgaliojimus; įsipareigojimą laikytis reikalavimų ir nuolat
tobulėti; roles ir atsakomybes; pagrindinius atitikties procesus; pranešimą apie pažeidimus; ir
neatitikties pasekmes. Aplink politiką išsidėsto kurso mokoma dokumentų hierarchija ---
pirmajame lygmenyje pati politika, antrajame teminės politikos, trečiajame procedūros --- ir
produktas kiekvieną įmonės dokumentą priskiria lygmeniui bei išvardija trūkstamus antrojo ir
trečiojo lygmens dokumentus, kiekvieną --- su teisiniu pagrindu, jei to dokumento reikalauja
kuri nors nuostata. Kodas neleidžia patvirtinti politikos, kol kuri nors iš aštuonių
dalių turi mažiau nei 150 simbolių arba joje nurodyto straipsnio nėra atitinkamame teisės akte.

\paragraph{2 kurso užduotis: valdymo struktūra, procesas ir atsakomybių matrica.} Antroji
užduoties pusė yra grandinė: dokumentuose įvardytos pareigybės, trūkstamos rolės, struktūra,
procesas, matrica. Produktas šią grandinę įgyvendina viename programos dokumente ir kodu
užtikrina pačios užduoties taisyklę matricai: \emph{kiekvienas etapas turi lygiai vieną
atskaitingą rolę ir bent vieną atsakingą rolę}. Jis taip pat tikrina, ar struktūra atitinka
kurso mokomas tris gynybos linijas --- antroji linija negali būti pavaldi pirmajai, o vidaus
auditas atsiskaito tik valdymo organui, --- reikalauja, kad procesas turėtų bent penkis etapus,
ir patikrina ISO~37301 punktus bei Lietuvos banko atitikties funkcijos gairių punktus, kuriais,
kaip nurodyta etapuose, jie remiasi.

\paragraph{3 kurso užduotis: rizikos registras.} Kursas kriterijus nustato tiksliai, ir
produkto konstantos yra kurso konstantos: tikimybė penkiais lygiais pagal keturis aspektus
(vidaus kontrolės būklė, teisės aiškumas ir priežiūros institucijos dėmesys, sandorio šalys,
dažnumas), poveikis penkiais lygiais pagal keturis aspektus (finansai, priežiūra ir
licencija, reputacija, mastas), balas kaip jų sandauga ir penki rizikos lygiai --- nuo 1 iki 3
nereikšminga, nuo 4 iki 6 toleruotina, nuo 8 iki 12 vidutinė, 15 ir 16 didelė, nuo 20 iki 25
kritinė, --- kiekvienas su jam leidžiamais atsako būdais. Scenarijus rašomas kurso sakiniu
\enquote{Tikėtinas [įvykis], kylantis iš [šaltinis], gali sukelti [pasekmės], kurios paveiks
[tikslas]} ir priskiriamas vienai iš šešių ISO~31022 A~priedo teisinės rizikos rūšių; produktas
reikalauja bent penkių scenarijų, apimančių bent keturias iš šešių rūšių, kiekvieną susietą su
įmonės padaliniu ir politika. Būtent skaičiuose produktas nesiremia vien modelio atsakymu:
modelis pateikia aštuonis lygius ir jų pagrindimą, o tikimybę, poveikį, balą ir rizikos lygį
apskaičiuoja kodas (\cref{ch:programme}). Registras eksportuojamas kurso nustatytais devyniais
stulpeliais.

\begin{humanbox}{Ko produktas už studentą nedaro}
Jis netvirtina savo paties darbo. Kiekvienas iš trijų programos dokumentų pateikiamas kaip
projektas, šalia išvardijant taisykles, kurias jis atitinka, ir taisykles, kurias pažeidžia;
redaguoja žmogus ir tvirtina žmogus, o PDF viršelyje nurodyta, ar kas nors tai padarė. Todėl iš
produkto pateikta kurso užduotis yra peržiūrėtas arba neperžiūrėtas projektas, ir joje pasakyta,
kuris iš jų.
\end{humanbox}

\section{Kurso sąvokos sistemoje}\label{sec:course:concepts}

Be užduočių, didžiąją sistemos formos dalį lėmė kurso paskaitos apie stebėseną ir technologo
rolę. \Cref{tab:course:concepts} išvardija sąvokas ir jas įkūnijančias sistemos dalis.

\begin{xltabular}{\textwidth}{@{}L{40mm}YL{20mm}@{}}
\caption{Kurso mokomos sąvokos ir kiekvieną jų įkūnijanti sistemos dalis}\label{tab:course:concepts}\\
\toprule
\hd{Kursas moko} & \hd{Sistema daro} & \hd{Skyrius} \\
\midrule
\endfirsthead
\toprule
\hd{Kursas moko} & \hd{Sistema daro} & \hd{Skyrius} \\
\midrule
\endhead
\bottomrule
\endlastfoot
Išvadą sudaro keturios dalys: būklė, kriterijus, priežastis ir pasekmė. &
  Ataskaitos išvada turi visas keturias. Būklė yra vertinimas ir jo būsena; kriterijus ---
  pareiga ir ją pagrindžianti nuostata; pasekmę kodas prijungia iš paties teisės akto sankcijų
  nuostatų; priežastis yra laukas, kurį pildo tik žmogus, nes dokumentai rodo, kad kažko
  trūksta, bet niekada --- kodėl. & \cref{ch:report} \\
Ataskaita atsako į tris klausimus: kas, kas iš to, ką daryti. &
  Išvados sako kas; pasekmė ir pareigos reikšmingumas sako, kas iš to; veiksmų planas, apimantis
  kiekvieną atvirą išvadą, sako ką daryti. PDF dokumentas išdėstytas nuo apibendrinimo žemyn. &
  \cref{ch:report} \\
Žmogus patvirtina darbo apimtį ir kiekvieno uždavinio rezultatą. &
  Ataskaita sustoja nusprendusi, kurie teisės aktai taikomi, ir laukia, kol žmogus patvirtins
  sąrašą, prieš pradėdama brangiąją dalį (žmogus gali leisti jai vykti iki galo, ir tada
  ataskaitoje įrašoma, kad niekas jos apimties nepatvirtino); žmogus peržiūri išvadas ir
  patvirtina baigtą ataskaitą. &
  \cref{ch:report} \\
Išvados ir veiksmai saugomi registre su būsenomis atviras, vykdomas, vėluojantis ir uždarytas,
  o veiksmas uždaromas tik po pakartotinio patikrinimo. &
  Veiksmų registras. Atsakingas asmuo gali pažymėti veiksmą atliktu, bet uždaro jį tik vėlesnė
  patvirtinta ataskaita --- pakartotinis patikrinimas, --- kurioje jo pareigos įvertintos kaip atitinkančios, nebent
  žmogus uždaro jį rankiniu būdu su pastaba, paaiškinančia kodėl. Vėlavimas nustatomas pagal
  terminą, rankiniu būdu jis niekada nežymimas. &
  \cref{ch:register} \\
Kriterijai yra ne tik teisė, bet ir pačios įmonės taisyklės. &
  Ataskaita gali vertinti, ar įmonės procedūros įgyvendina tai, ko reikalauja jos pačios
  politikos, kaip atskirą išvadų grupę, skaičiuojamą atskirai nuo teisinių. & \cref{ch:report} \\
Trys gynybos linijos. &
  Valdymo struktūros programos dokumentas kiekvieną organą, funkciją ir padalinį priskiria
  linijai ir neleidžia patvirtinti struktūros, kurioje antroji linija pavaldi pirmajai. &
  \cref{ch:programme} \\
Rizika yra tikimybės ir poveikio sandauga pagal nustatytus kriterijus. &
  Rizikų programos dokumentas, kuriame balą ir rizikos lygį apskaičiuoja kodas. & \cref{ch:programme} \\
\end{xltabular}

\section{Penkios technologo pareigos}\label{sec:course:duties}

Pirmojoje kurso paskaitoje DI ir atitikties sistemų teisinių operacijų technologas apibūdinamas
kaip asmuo, kuris kuria, testuoja ir diegia instrukcijomis grįstus tikrinimo procesus ---
procesus, kurie tikrina įmonės dokumentus pagal jos vidaus politikas ir pagal teisę, --- ir
įvardijamos penkios jo pareigos. Tai tiesiausiai išreiškia, ko kursas tikisi iš tokios
sistemos kaip ši. \Cref{tab:course:duties} sistemą palygina su kiekviena pareiga, nurodydama ir
tai, kur ji nepasiekia tikslo.

\begin{xltabular}{\textwidth}{@{}L{34mm}YY@{}}
\caption{Penkios paskaitoje įvardytos technologo pareigos ir sistema}\label{tab:course:duties}\\
\toprule
\hd{Pareiga, kaip ją nusako paskaita} & \hd{Ką daro sistema} & \hd{Ko sistema nepadaro} \\
\midrule
\endfirsthead
\toprule
\hd{Pareiga, kaip ją nusako paskaita} & \hd{Ką daro sistema} & \hd{Ko sistema nepadaro} \\
\midrule
\endhead
\bottomrule
\endlastfoot
\textbf{Instrukcijų inžinerija ir schemų laikymosi užtikrinimas:} tikslios, struktūrizuotos, JSON
  schemomis apribotos instrukcijos, kurios tikrina įmonės dokumentus ir pažymi neatitinkančias
  nuostatas ar faktų nenuoseklumus. &
  Dvidešimt viena iš 23 instrukcijų atsako pagal griežtą JSON schemą, kurios laikytis prašomas
  tiekėjas ir kurią asistento paslauga dar kartą patikrina (\cref{ch:models,ch:prompts}).
  Atitikties spragų analizė kiekvienai pareigai priskiria būseną ir pažodžiui cituoja rastą
  įrodymą; dokumento inventorius išvardija jo silpnąsias vietas ir nuorodas į teisę; nuorodų
  patikra pažymi nuorodą į straipsnį, kuriame nėra to, ką teigia dokumentas
  (\cref{ch:report,ch:company}). &
  Dvi pokalbio instrukcijos atsako laisvu tekstu. \\
\textbf{Pagrindimas šaltiniais ir konteksto tikrinimas:} paieška papildytos instrukcijos,
  kurios dokumentus tiesiogiai lygina su vidaus politikų rinkiniais, standartinėmis veiklos
  procedūromis ir taikytinais teisės aktais. &
  Teisė imama iš teisės aktų bazės, ne iš modelio atminties, o nuostata, kurią modelis
  nurodo iš atminties, --- programos dokumento teisinis pagrindas --- ten patikrinama: pokalbis
  nuostatas suranda mišriąja paieška; ataskaita modeliui perduoda pačias nuostatas arba iš jų
  išskirtas pareigas (\cref{ch:corpus}). Su kiekvienu atitikties spragų analizės iškvietimu
  siunčiamas visas įmonės dokumentų rinkinys --- iš pradžių buvo naudojamos surastos ištraukos,
  bet jų atsisakyta, nes jos praleisdavo dokumentuose esančius įrodymus (\cref{ch:report}).
  Kriterijai gali būti ir pačios įmonės politikos. &
  Paieška lieka atsarginiu sprendimu dokumentų rinkiniui, kuris per didelis, kad būtų siunčiamas visas, o
  pokalbis ir toliau dirba su tuo, ką suranda. \\
\textbf{Automatinis vertinimas ir apsauga nuo haliucinacijų:} regresinių testų rinkiniai ---
  modelis kaip vertintojas, natūraliosios kalbos išvada (angl. \en{natural-language
  inference}), vienetų testai, --- kurie nuolat matuoja instrukcijų patikimumą, nuorodų
  tikslumą ir kraštutinių atvejų tvarkymą. &
  Šeši atvejų rinkiniai, kurių atvejus 2026 m. lapkričio 12 d. pažymėjo DI modelis ir kurių dar
  nepatvirtino joks žmogus, vykdomi su tikruoju modeliu per paties produkto funkcijas ir
  lyginami su įrašytu baziniu rezultatu; trijose vietose vertina antrasis modelio
  iškvietimas; apie 6\,100 testų vykdoma nekviečiant modelio (\cref{ch:control,ch:testing}).
  Apsauga, nepriklausanti nuo jokio modelio: kiekviena citata, kuria remiasi išvada ar
  programos dokumentas, surandama jos dokumente, kiekviena struktūrizuotame atsakyme nurodyta
  nuostata patikrinama, kiekvienas balas apskaičiuojamas. &
  Natūraliosios kalbos išvados modelis nenaudojamas. Rinkiniai matuoja šešias iš 23
  instrukcijų ir vykdomi rankiniu būdu prieš keičiant instrukciją ar modelį ir po to, o ne
  nuolat. \\
\textbf{Bendradarbiavimas su ekspertais ir etaloniniai rinkiniai:} sudėtingų reguliavimo
  reikalavimų pavertimas deterministinėmis instrukcijomis kartu su atitikties pareigūnais ir
  teisininkais bei patikrintų etaloninių duomenų rinkinių saugojimas. &
  Pareigų registras yra reikalavimai deterministine forma: kiekviena pareiga užrašyta vieną
  kartą, susieta su savo nuostata, įvertinta ir peržiūrėta bei pakartotinai naudojama kiekvienoje
  ataskaitoje (\cref{ch:obligations}). Kiekviena būsena ir verdiktas, kuriuos gali pateikti
  instrukcija, apibrėžti pačioje instrukcijoje (\cref{ch:prompts}). Pažymėti atvejai saugomi
  repozitorijoje, penkiuose iš šešių rinkinių --- su kiekvienos žymos priežastimi. &
  Atvejus 2026 m. lapkričio 12 d. pažymėjo DI programavimo asistentas, kuriuo sistema buvo
  kuriama (Claude), skaitydamas šaltinius; kartu pateikiamas registras peržiūrėtas mašininiu
  būdu, įvertintas vertintojo ir žmogaus neperskaitytas įrašas po įrašo
  (\cref{ch:obligations,ch:control}). Nei žymų, nei registro dar nepatvirtino joks atitikties
  pareigūnas ar teisininkas. Abu parengti taip, kad žmogus galėtų pakeisti bet kurią žymą ar
  įrašą. \\
\textbf{Audituojamumas ir integracija:} tikrinimas, įdiegtas į esamas teisinių operacijų
  sistemas --- dokumentų valdymo, sutarčių valdymo sistemas, procesų variklius, --- su visišku įvykių
  registravimu ir audito sekomis. &
  Kiekvienas bet kurio asmens pakeitimas yra veiksmų žurnale; kiekvienas atsakęs modelio
  iškvietimas registruojamas su instrukcijos versija, žetonais ir kaina (išskyrus paieškos
  užklausos vektorizavimą), o kiekvienas pokalbio ėjimas --- su stenograma; kiekviena
  ataskaita įrašo instrukcijų ir teisės aktų redakcijas, iš kurių ji parengta
  (\cref{ch:security,ch:models}). Visa sistema pasiekiama per JSON API (\cref{ch:api}). &
  Nesukurta jokia jungtis su dokumentų ar sutarčių valdymo sistema: dokumentai įkeliami. \\
\end{xltabular}

\section{Kur sistema seka paties šaltinio tekstą}\label{sec:course:sources}

Dviem atvejais sistemos elgsena kyla iš paties šaltinio formuluotės, ir skaitytojas, lyginantis
ją su to šaltinio santrauka, gali tai pastebėti.

\begin{itemize}
  \item \textbf{Lietuvos banko atitikties funkcijos gairės} suformuluotos kaip rekomendacijos
    (\enquote{turėtų}) ir skirtos jų trečiajame punkte išvardytų rūšių įstaigoms, tarp kurių
    nėra kriptoturto paslaugų teikėjų. Kataloge jų apimtis nustatyta atitinkamai. Todėl tokiai
    įmonei kaip \demo{} ataskaita jų savaime nevertina; žmogus apimties ekrane gali įtraukti
    jas į ataskaitos apimtį, ir tada jos vertinamos kaip geroji praktika ir taip pažymimos
    (\cref{ch:report}).
  \item \textbf{ISO~37301 punktų numeriai.} Atitikties proceso etapai nurodo standarto punktus
    taip, kaip juos numeruoja pats standartas: susirūpinimą keliančių klausimų iškėlimas yra
    8.3~punktas, o tyrimo procesai --- 8.4~punktas. Etapą, kuris nenurodo nei standarte
    esančio punkto, nei gairių punkto, paties produkto patikra pažymi įspėjimu
    (\cref{ch:programme}).
\end{itemize}

\section{Ką sistema prideda prie kurso metodo}\label{sec:course:adds}

Keturių sistemos dalykų neprašo jokia užduotis; jie kyla iš to, kad metodas rimtai įgyvendintas
kaip programinė įranga.

\begin{description}
  \item[Teisė kaip prižiūrima teisės aktų bazė.] Užduotys leidžia modeliui kalbėti apie teisę.
    Sistema atsisiunčia teisės aktus iš oficialių šaltinių, saugo jų redakcijas ir įrašo, ką
    pakeitė kiekviena nauja redakcija, --- kad galėtų pasakyti, kurios išvados remiasi nuo to
    laiko pasikeitusia nuostata (\cref{ch:corpus,ch:register}).
  \item[Pareigų registras.] Tarp teisės ir ataskaitos yra peržiūrėtas sąrašas, ką nustato
    kiekvienas teisės aktas, išskirtas vieną kartą ir pakartotinai naudojamas kiekvienoje
    ataskaitoje. Būtent jis daro dviejų įmonių ataskaitas palyginamas ir suteikia žmogui vieną
    vietą, kur pataisyti teisės skaitymą (\cref{ch:obligations}).
  \item[Teisės aktų nuorodų patikra.] Pavyzdinės organizacijos apžvalginis dokumentas kai kuriuos
    reglamento dėl kriptoturto rinkų straipsnius nurodo neteisingais numeriais. Sistema
    kiekvieną dokumente nurodytą straipsnį patikrina pagal teisės aktų bazę ir pagal tai, ką,
    dokumento teigimu, tas straipsnis apima (\cref{ch:company}).
  \item[Kaina kaip projektavimo apribojimas.] Kiekvienas sistemos įrašomas iškvietimas
    įkainojamas; ataskaita prieš pradėdama parodo savo kainos įvertį; brangioji dalis laukia
    žmogaus, nebent vykdymas pradėtas be apimties peržiūros; vienkartinis darbas su teise
    atskirtas nuo darbo kiekvienai įmonei (\cref{ch:models,ch:control}).
\end{description}

\section{Ko kursas moko, bet sistema nedaro}\label{sec:course:not}

\begin{itemize}
  \item \textbf{Metinė atitikties ataskaita ir stebėsenos programa.} Kursas abi laiko atitikties
    funkcijos produktais. Sistema turi joms reikalingą medžiagą --- ataskaitas, veiksmų
    registrą, kiekvienos pareigos istoriją, --- bet nė vienos jų nerengia.
  \item \textbf{Atitikties valdymo sistemos brandos vertinimas} nėra sistemos produktas.
  \item \textbf{Patys ISO standartai.} Sistema gali suskaidyti standartą į punktus ir išskirti
    jo reikalavimus (\cref{ch:corpus,ch:obligations}), o įmonė gali jį taikyti ir būti pagal jį
    vertinama. Kurso turimos ISO standartų kopijos yra nuskenuoti vaizdai, naudojami pagal
    mokymo licenciją: sistema neskaito skenuotų dokumentų, o licencija neapima jų įkėlimo.
    Įdiegtoje teisės aktų bazėje nėra nė vieno standarto.
  \item \textbf{Veikimo efektyvumas.} Kursas skiria kontrolės priemonės dizainą nuo jos
    veikimo. Ataskaita vertina dizainą pagal dokumentus; ji pažymi, kur reikėtų patikrinti
    veikimą, ir patį patikrinimą palieka žmogui.
  \item \textbf{Išvados priežastis.} Sąmoningai palikta peržiūrėtojui.
\end{itemize}

Pačių DI kelių apribojimai --- kas nematuojama, ko vertintojas negali nustatyti, kaip buvo
peržiūrėtas kartu pateikiamas pareigų registras --- surinkti \cref{ch:control}.
